% One-page project summary, BMEN 600 Group 11. Sent as an attachment; the request itself is in the email (email.md).
% Numbers: run 20261113_124400 of the group repository (main, commit 2585c0b), printed by export_numbers.py.
% Method details: config.json of that run and src/models.py, src/training.py, src/ann_to_snn.py at that commit.
% No reference list on the page. The sources of the Motivation claims are in ../references/references_fixed.md
% ([8] and [9] conversion, [11] speech and gesture, [12] and [13] motor cortex spike trains, [15] the study on our dataset).
% Build:  latexmk -pdf summary.tex     (must stay on one page: pdfinfo summary.pdf)
\documentclass[10pt,letterpaper]{article}

\usepackage[letterpaper,top=0.45in,bottom=0.45in,left=0.7in,right=0.7in]{geometry}
\usepackage[T1]{fontenc}
\usepackage{mathptmx}            % Times text and maths, as in most journals
\usepackage{microtype}
\usepackage{xcolor}
\usepackage{booktabs}
\usepackage{array}
\usepackage[font=footnotesize,labelfont=bf,labelsep=period,justification=justified,singlelinecheck=false,skip=3pt]{caption}
\usepackage{tikz}
\usepackage{pgfplots}
\usepackage{pgfplotstable}
\usepackage[hidelinks]{hyperref}
\pgfplotsset{compat=1.18}

% colour-blind-safe categorical slots 1 to 3 (figure markers only)
\definecolor{snncol}{HTML}{2A78D6}
\definecolor{rsnncol}{HTML}{EB6834}
\definecolor{anncol}{HTML}{1BAF7A}
\definecolor{rangecol}{HTML}{C4CAD3}

\setlength{\parindent}{0pt}
\setlength{\parskip}{3pt}
\pagestyle{empty}
\newcommand{\head}[1]{\textbf{#1.}}
\newcommand{\sect}[1]{\par\vspace{3pt}{\large\bfseries #1}\par\vspace{-1pt}}

\begin{document}

\fontsize{9.7}{11.6}\selectfont
{\Large\bfseries Recurrent versus converted spiking neural networks for motor imagery EEG decoding}\\[2pt]
{\bfseries A parameter-matched, cross-session comparison on 4 classes}
\hfill {\small Project summary, 18 November 2026. Preliminary results.}\\[1pt]
Lex Berezowski, Amirsaeed Ghodsinejad and Pahansith Kiriwaththuduwage \hfill BMEN 600, Group 11, University of Calgary
\par\vspace{1pt}\hrule

\sect{Motivation}
Spiking neural networks (SNNs) communicate with binary spikes and can be energy efficient on neuromorphic hardware,
so they are candidates for decoding in wearable brain--computer interfaces (BCIs). For motor imagery from the electroencephalogram (EEG), most SNNs are feedforward.
Many are made by ANN-to-SNN conversion of a trained convolutional network, and the accuracy is usually about the same.
Recurrent SNNs (RSNNs) trained with surrogate gradients achieve good results on time-domain tasks such as speech and
gesture recognition, and on spike trains from the primary motor cortex of monkeys. On those spike trains, RSNNs were
more accurate than feedforward SNNs. For EEG, we have very limited evidence. On our dataset, one study compared a
recurrent model with a convolutional SNN. But it used k-fold cross-validation within each subject, so the training
trials and the test trials came from the same recordings. It also did not test if the difference between the models
was statistically significant. A real BCI is trained on one day and used on another day. To our knowledge, no study
has tested if an RSNN is better than a feedforward SNN in that case.

\head{Research question} With an equal number of parameters, and with training and test on sessions recorded on
different days (cross-session evaluation), does a directly trained RSNN decode 4-class motor imagery more accurately
than a feedforward SNN converted from EEGNet?

\sect{Methods}
\head{Data} BCI Competition IV dataset 2a: 9 subjects, 22 EEG channels, 4 motor imagery tasks (left hand, right hand,
both feet, tongue), 2 sessions of 288 trials per subject, recorded on different days. Signals were band-pass filtered
(4--40\,Hz) and resampled to 128\,Hz. Each trial is the 2\,s window from 0.5\,s after cue onset
(22\,$\times$\,256 samples). Trials marked as artifacts were kept.

\head{Evaluation protocol} For each target subject, a model was pretrained on session~1 of 5 of the 8 other subjects
(the 3 others gave the validation set, up to 500 epochs), then fine-tuned on session~1 of the target subject
(80\,\% training, 20\,\% validation, 100 epochs), then tested once on session~2 of the target subject (288 trials). This
was repeated for 10 random choices of the pretraining subjects: 90 folds per model, the same folds for all models.
Training used Adam, cross-entropy loss, batch size 16, learning rate 0.001 and gradient clipping. The weights
with the lowest validation loss were kept. Session~2 was never used for training or model selection.

\head{Models} \textbf{(a)~EEGNet-4,2}: a compact convolutional network with depthwise and separable
convolutions, 796 parameters. \textbf{(b)~Converted SNN}: ANN-to-SNN conversion of the fine-tuned EEGNet of the same
fold. Batch normalization was integrated into the weights, the weights were normalized with the 99.9th percentile of
the activations, and the units were replaced by integrate-and-fire neurons, simulated for 256 time steps; 776
parameters, no further training. \textbf{(c)~RSNN}: a fully connected input layer, one fully connected recurrent layer
of 17 leaky integrate-and-fire (LIF) neurons, and a linear readout of the firing rates; 752 parameters, matched
to (b); trained directly by backpropagation through time with a surrogate gradient. \textbf{(d)~RNN control}: the same
layers as (c) with non-spiking units, 752 parameters. \textbf{(e)~Neuron-matched RSNN}: as (c) with 2,560 LIF
neurons, the same number of neurons as EEGNet (6.6 million parameters).

\head{Statistics} The subject is the unit of analysis: the 10 repeats of a subject were averaged, which gives 9 paired
values per model. Models were compared with the exact two-sided Wilcoxon signed-rank test, with Holm correction over the
three planned comparisons. An accuracy is significantly above chance if it is above 0.292, the threshold from the
binomial distribution (288 trials, 4 classes, $\alpha = 0.05$).

\sect{Results}
\begin{minipage}[t]{0.53\linewidth}\vspace{0pt}
\captionof{table}{Performance on the test session (mean $\pm$ standard deviation over 9 subjects).
$\kappa$: Cohen's kappa. Last column: subjects above the chance threshold.}
\footnotesize\setlength{\tabcolsep}{3.5pt}
\begin{tabular}{@{}lrccc@{}}
\toprule
Model & Parameters & Accuracy & $\kappa$ & Above chance \\
\midrule
(a) EEGNet                & 796       & 0.600 $\pm$ 0.150 & 0.466 & 9 of 9 \\
(b) Converted SNN         & 776       & 0.565 $\pm$ 0.154 & 0.420 & 9 of 9 \\
(c) RSNN                  & 752       & 0.385 $\pm$ 0.100 & 0.180 & 6 of 9 \\
(d) RNN control           & 752       & 0.411 $\pm$ 0.106 & 0.215 & 7 of 9 \\
(e) Neuron-matched RSNN   & 6,622,724 & 0.341 $\pm$ 0.085 & 0.122 & 5 of 9 \\
\bottomrule
\end{tabular}

\vspace{5pt}\fontsize{9.7}{11.6}\selectfont
The converted SNN kept most of the accuracy of EEGNet (difference 0.034, $p = 0.098$). The RSNN was less accurate than
the converted SNN for all 9 subjects (difference 0.180, Holm-corrected $p = 0.012$, the smallest value possible with
9 subjects). The non-spiking RNN control was only slightly above the RSNN (difference 0.026, uncorrected $p = 0.074$).
More neurons did not help: the neuron-matched RSNN was below the 17-neuron RSNN for all 9 subjects.
\end{minipage}\hfill
\begin{minipage}[t]{0.44\linewidth}\vspace{0pt}
\pgfplotstableread[col sep=comma]{
row,snn,rsnn,ann,wide,mid,half
1,0.618,0.453,0.712,0.368,0.5403,0.1719
2,0.400,0.281,0.409,0.263,0.3363,0.0731
3,0.734,0.477,0.781,0.419,0.6000,0.1812
4,0.463,0.374,0.524,0.333,0.4281,0.0955
5,0.366,0.275,0.413,0.262,0.3373,0.0755
6,0.416,0.284,0.458,0.251,0.3547,0.1033
7,0.679,0.314,0.628,0.273,0.4758,0.2029
8,0.656,0.515,0.737,0.423,0.5797,0.1571
9,0.756,0.493,0.736,0.480,0.6182,0.1380
10,0.565,0.385,0.600,0.341,0.4706,0.1293
}\persubject
\centering
\begin{tikzpicture}
\begin{axis}[
  width=\linewidth, height=5.0cm,
  font=\scriptsize,
  xmin=0.2, xmax=0.8, ymin=0.4, ymax=10.6, y dir=reverse,
  ytick={1,...,10},
  yticklabels={Subject 1,Subject 2,Subject 3,Subject 4,Subject 5,Subject 6,Subject 7,Subject 8,Subject 9,\textbf{Mean}},
  xtick={0.2,0.3,...,0.8},
  xticklabel style={/pgf/number format/fixed,/pgf/number format/precision=1},
  xlabel={Test accuracy (session 2)}, xlabel style={yshift=2pt},
  xmajorgrids, major grid style={color=black!10},
  axis line style={color=black!50}, tick style={color=black!50},
  legend style={at={(1,1.02)}, anchor=south east, draw=none, legend columns=2, font=\scriptsize, inner sep=0pt,
                /tikz/every even column/.append style={column sep=6pt}},
  legend cell align=left,
]
  \draw[dashed, color=black!60, line width=0.6pt] (axis cs:0.2917,0.4) -- (axis cs:0.2917,10.6);
  \draw[color=black!40] (axis cs:0.2,9.5) -- (axis cs:0.8,9.5);
  \addplot[forget plot, draw=none, mark=none,
           error bars/.cd, x dir=both, x explicit, error mark=none,
           error bar style={line width=1.6pt, color=rangecol}]
    table[x=mid, y=row, x error=half] {\persubject};
  \addplot[only marks, mark=*, mark size=2.1pt, mark options={fill=anncol, draw=white, line width=0.6pt}]
    table[x=ann, y=row] {\persubject};
  \addlegendentry{(a) EEGNet}
  \addplot[only marks, mark=*, mark size=2.1pt, mark options={fill=snncol, draw=white, line width=0.6pt}]
    table[x=snn, y=row] {\persubject};
  \addlegendentry{(b) Converted SNN}
  \addplot[only marks, mark=*, mark size=2.1pt, mark options={fill=rsnncol, draw=white, line width=0.6pt}]
    table[x=rsnn, y=row] {\persubject};
  \addlegendentry{(c) RSNN}
  \addplot[only marks, mark=o, mark size=1.9pt, mark options={draw=rsnncol, line width=0.9pt}]
    table[x=wide, y=row] {\persubject};
  \addlegendentry{(e) Neuron-matched RSNN}
\end{axis}
\end{tikzpicture}
\captionof{figure}{Test accuracy of each subject (mean of 10 repeats). Dashed line: chance threshold 0.292.}
\end{minipage}

\sect{Interpretation and limits}
ANN-to-SNN conversion of a compact convolutional network gave a much better spiking decoder than direct training of
a one-layer RSNN, and neither more neurons nor non-spiking units closed the gap. This suggests that the limit is the single recurrent layer applied directly to the EEG samples, and not spiking.
The models also differ in architecture and in
training method, so the difference cannot be attributed to recurrence alone. The RSNN settings came from our research
plan and were not tuned. There are 9 subjects from one offline dataset. Results are preliminary until our pipeline checks are complete.

\end{document}
